\small
\begin{enumerate}[label=\textbf{\arabic*.}, leftmargin=0pt, itemsep=1.2em]
\item Что в генетическом алгоритме называют хромосомой?
\begin{itemize}[label=$\square$, leftmargin=1.5cm, itemsep=0.2em, topsep=0.3em]
\item Кодированное представление кандидата на решение задачи
\item Только целевую функцию оптимизации
\item Массив значений ошибки на всех эпохах
\item Матрицу весов нейронной сети
\end{itemize}
{\small\textit{Правильный ответ: Кодированное представление кандидата на решение задачи.}}

\item Для чего в классическом ГА используется бинарное кодирование вещественной переменной?
\begin{itemize}[label=$\square$, leftmargin=1.5cm, itemsep=0.2em, topsep=0.3em]
\item Для представления значения на отрезке в виде битовой строки
\item Для преобразования задачи классификации в задачу кластеризации
\item Для вычисления градиента функции
\item Для построения матрицы путаницы
\end{itemize}
{\small\textit{Правильный ответ: Для представления значения на отрезке в виде битовой строки.}}

\item От чего зависит точность представления вещественного значения при бинарном кодировании?
\begin{itemize}[label=$\square$, leftmargin=1.5cm, itemsep=0.2em, topsep=0.3em]
\item От длины хромосомы в битах
\item Только от числа поколений
\item Только от размера популяции
\item Только от вероятности мутации
\end{itemize}
{\small\textit{Правильный ответ: От длины хромосомы в битах.}}

\item В чем состоит идея пропорциональной селекции (метода рулетки)?
\begin{itemize}[label=$\square$, leftmargin=1.5cm, itemsep=0.2em, topsep=0.3em]
\item Вероятность выбора особи пропорциональна ее приспособленности
\item Всегда выбирается только худшая особь
\item Выбор родителей не зависит от фитнеса
\item Каждая особь выбирается ровно один раз
\end{itemize}
{\small\textit{Правильный ответ: Вероятность выбора особи пропорциональна ее приспособленности.}}

\item Что делает одноточечный кроссинговер?
\begin{itemize}[label=$\square$, leftmargin=1.5cm, itemsep=0.2em, topsep=0.3em]
\item Обменивает хвостовые части двух родительских хромосом после выбранной точки
\item Инвертирует каждый бит хромосомы
\item Удаляет половину популяции
\item Случайно меняет местами двух родителей
\end{itemize}
{\small\textit{Правильный ответ: Обменивает хвостовые части двух родительских хромосом после выбранной точки.}}

\item Как работает побитовая мутация в бинарном ГА?
\begin{itemize}[label=$\square$, leftmargin=1.5cm, itemsep=0.2em, topsep=0.3em]
\item Инвертирует отдельные биты с малой вероятностью
\item Удаляет худшие особи из популяции
\item Копирует лучшую особь без изменений
\item Усредняет две хромосомы поэлементно
\end{itemize}
{\small\textit{Правильный ответ: Инвертирует отдельные биты с малой вероятностью.}}

\item Зачем в генетическом алгоритме нужна мутация?
\begin{itemize}[label=$\square$, leftmargin=1.5cm, itemsep=0.2em, topsep=0.3em]
\item Для поддержания разнообразия и выхода из локальных оптимумов
\item Чтобы полностью заменить кроссинговер
\item Чтобы сделать все особи одинаковыми
\item Чтобы уменьшить размер хромосомы
\end{itemize}
{\small\textit{Правильный ответ: Для поддержания разнообразия и выхода из локальных оптимумов.}}

\item Что означает термин "поколение" в генетическом алгоритме?
\begin{itemize}[label=$\square$, leftmargin=1.5cm, itemsep=0.2em, topsep=0.3em]
\item Одна итерация эволюционного обновления популяции
\item Одна операция мутации над битом
\item Один элемент хромосомы
\item Одно значение целевой функции
\end{itemize}
{\small\textit{Правильный ответ: Одна итерация эволюционного обновления популяции.}}

\item Как обычно интерпретируется фитнес-функция при задаче максимизации?
\begin{itemize}[label=$\square$, leftmargin=1.5cm, itemsep=0.2em, topsep=0.3em]
\item Чем больше значение фитнеса, тем лучше особь
\item Чем меньше значение фитнеса, тем лучше особь
\item Фитнес не связан с качеством решения
\item Фитнес равен длине хромосомы
\end{itemize}
{\small\textit{Правильный ответ: Чем больше значение фитнеса, тем лучше особь.}}

\item Зачем используется элитарная стратегия?
\begin{itemize}[label=$\square$, leftmargin=1.5cm, itemsep=0.2em, topsep=0.3em]
\item Чтобы гарантированно переносить лучшие найденные решения в следующее поколение
\item Чтобы сделать все особи случайными
\item Чтобы исключить мутацию из алгоритма
\item Чтобы запретить отбор родителей
\end{itemize}
{\small\textit{Правильный ответ: Чтобы гарантированно переносить лучшие найденные решения в следующее поколение.}}

\item Что обычно происходит при слишком малой вероятности мутации?
\begin{itemize}[label=$\square$, leftmargin=1.5cm, itemsep=0.2em, topsep=0.3em]
\item Растет риск преждевременной сходимости
\item Алгоритм превращается в случайный поиск
\item Хромосомы перестают участвовать в кроссинговере
\item Размер популяции автоматически увеличивается
\end{itemize}
{\small\textit{Правильный ответ: Растет риск преждевременной сходимости.}}

\item Что обычно происходит при слишком большой вероятности мутации?
\begin{itemize}[label=$\square$, leftmargin=1.5cm, itemsep=0.2em, topsep=0.3em]
\item Поиск становится близким к случайному блужданию
\item Решение обязательно становится оптимальным
\item Элитизм перестает работать математически
\item Хромосомы становятся вещественными
\end{itemize}
{\small\textit{Правильный ответ: Поиск становится близким к случайному блужданию.}}

\item Как увеличение длины бинарной хромосомы влияет на поиск?
\begin{itemize}[label=$\square$, leftmargin=1.5cm, itemsep=0.2em, topsep=0.3em]
\item Повышает точность представления, но усложняет пространство поиска
\item Уменьшает точность и сокращает число локальных экстремумов
\item Отменяет необходимость кроссинговера
\item Полностью устраняет риск преждевременной сходимости
\end{itemize}
{\small\textit{Правильный ответ: Повышает точность представления, но усложняет пространство поиска.}}

\item Что называют генотипом в терминах ГА?
\begin{itemize}[label=$\square$, leftmargin=1.5cm, itemsep=0.2em, topsep=0.3em]
\item Кодовую структуру решения, например битовую строку
\item График целевой функции
\item Значение лучшего фитнеса
\item Само оптимальное решение в вещественной форме без кодирования
\end{itemize}
{\small\textit{Правильный ответ: Кодовую структуру решения, например битовую строку.}}

\item Что называют фенотипом в терминах ГА?
\begin{itemize}[label=$\square$, leftmargin=1.5cm, itemsep=0.2em, topsep=0.3em]
\item Интерпретацию закодированной хромосомы в терминах решаемой задачи
\item Список всех родителей популяции
\item Вероятность кроссинговера
\item Номер поколения
\end{itemize}
{\small\textit{Правильный ответ: Интерпретацию закодированной хромосомы в терминах решаемой задачи.}}

\item Как влияет увеличение размера популяции?
\begin{itemize}[label=$\square$, leftmargin=1.5cm, itemsep=0.2em, topsep=0.3em]
\item Обычно улучшает покрытие пространства поиска и надежность поиска
\item Всегда ухудшает результат
\item Делает мутацию ненужной
\item Приводит к линейной разделимости данных
\end{itemize}
{\small\textit{Правильный ответ: Обычно улучшает покрытие пространства поиска и надежность поиска.}}

\item Что является типичным критерием остановки ГА?
\begin{itemize}[label=$\square$, leftmargin=1.5cm, itemsep=0.2em, topsep=0.3em]
\item Достижение заданного числа поколений или стабилизация качества решений
\item Зануление всех весов
\item Появление хотя бы одной мутации
\item Преобразование хромосом в вещественные числа
\end{itemize}
{\small\textit{Правильный ответ: Достижение заданного числа поколений или стабилизация качества решений.}}

\item Почему генетический алгоритм называют эволюционным поиском?
\begin{itemize}[label=$\square$, leftmargin=1.5cm, itemsep=0.2em, topsep=0.3em]
\item Потому что он моделирует отбор, скрещивание и мутацию популяции решений
\item Потому что он использует только производные функции
\item Потому что он работает только с биологическими данными
\item Потому что он сортирует решения по алфавиту
\end{itemize}
{\small\textit{Правильный ответ: Потому что он моделирует отбор, скрещивание и мутацию популяции решений.}}

\item Какое из утверждений о глобальном и локальном экстремуме верно?
\begin{itemize}[label=$\square$, leftmargin=1.5cm, itemsep=0.2em, topsep=0.3em]
\item ГА стремится избежать застревания в локальных экстремумах за счет разнообразия популяции
\item Локальный экстремум всегда совпадает с глобальным
\item Наличие локальных экстремумов не влияет на поиск
\item ГА может работать только с выпуклыми функциями
\end{itemize}
{\small\textit{Правильный ответ: ГА стремится избежать застревания в локальных экстремумах за счет разнообразия популяции.}}

\item Какова роль кроссинговера в ГА?
\begin{itemize}[label=$\square$, leftmargin=1.5cm, itemsep=0.2em, topsep=0.3em]
\item Комбинировать полезные фрагменты решений от разных родителей
\item Вычислять значение целевой функции
\item Нормализовать признаки выборки
\item Определять размер окна временного ряда
\end{itemize}
{\small\textit{Правильный ответ: Комбинировать полезные фрагменты решений от разных родителей.}}

\item Какой диапазон точек разреза допустим для одноточечного кроссинговера хромосомы длины L?
\begin{itemize}[label=$\square$, leftmargin=1.5cm, itemsep=0.2em, topsep=0.3em]
\item От 1 до L-1
\item Только 0
\item Только L
\item Только середина строки
\end{itemize}
{\small\textit{Правильный ответ: От 1 до L-1.}}

\item Что дает использование бинарного кодирования в задаче одномерной оптимизации?
\begin{itemize}[label=$\square$, leftmargin=1.5cm, itemsep=0.2em, topsep=0.3em]
\item Унифицированное дискретное представление пространства поиска
\item Гарантию нахождения точного глобального максимума
\item Отсутствие необходимости в функции приспособленности
\item Возможность обучать сеть без учителя
\end{itemize}
{\small\textit{Правильный ответ: Унифицированное дискретное представление пространства поиска.}}

\item Какой баланс должен поддерживаться в ГА?
\begin{itemize}[label=$\square$, leftmargin=1.5cm, itemsep=0.2em, topsep=0.3em]
\item Баланс между исследованием новых областей и использованием уже найденных хороших решений
\item Баланс между количеством входов и выходов нейрона
\item Баланс между числом классов и размером тестовой выборки
\item Баланс между средним и медианой данных
\end{itemize}
{\small\textit{Правильный ответ: Баланс между исследованием новых областей и использованием уже найденных хороших решений.}}

\item Почему сохранение лучшей особи полезно?
\begin{itemize}[label=$\square$, leftmargin=1.5cm, itemsep=0.2em, topsep=0.3em]
\item Потому что хорошие решения могут быть потеряны из-за случайности операторов
\item Потому что это исключает необходимость оценивать популяцию
\item Потому что после этого все потомки становятся оптимальными
\item Потому что это уменьшает длину хромосомы
\end{itemize}
{\small\textit{Правильный ответ: Потому что хорошие решения могут быть потеряны из-за случайности операторов.}}

\item Что происходит после декодирования бинарной строки?
\begin{itemize}[label=$\square$, leftmargin=1.5cm, itemsep=0.2em, topsep=0.3em]
\item Получается значение переменной в исходном диапазоне задачи
\item Хромосома удаляется из популяции
\item Задача становится задачей классификации
\item Происходит нормализация выборки
\end{itemize}
{\small\textit{Правильный ответ: Получается значение переменной в исходном диапазоне задачи.}}

\item Какое утверждение о ГА наиболее корректно?
\begin{itemize}[label=$\square$, leftmargin=1.5cm, itemsep=0.2em, topsep=0.3em]
\item ГА чаще всего находит хорошие квазиоптимальные решения, но не гарантирует точный оптимум
\item ГА всегда находит глобальный оптимум за фиксированное число шагов
\item ГА работает только с непрерывными функциями
\item ГА не зависит от выбора параметров
\end{itemize}
{\small\textit{Правильный ответ: ГА чаще всего находит хорошие квазиоптимальные решения, но не гарантирует точный оптимум.}}

\item Какой фактор сильнее всего влияет на давление отбора в методе рулетки?
\begin{itemize}[label=$\square$, leftmargin=1.5cm, itemsep=0.2em, topsep=0.3em]
\item Различие значений фитнеса между особями
\item Количество скрытых слоев сети
\item Число бит в одном гене
\item Форма графика функции активации
\end{itemize}
{\small\textit{Правильный ответ: Различие значений фитнеса между особями.}}

\item Что называют преждевременной сходимостью популяции?
\begin{itemize}[label=$\square$, leftmargin=1.5cm, itemsep=0.2em, topsep=0.3em]
\item Ситуацию, когда разнообразие популяции исчезает слишком рано и поиск замирает
\item Ситуацию, когда мутация становится равной нулю автоматически
\item Ситуацию, когда кроссинговер выполняется только один раз
\item Ситуацию, когда все особи становятся вещественными
\end{itemize}
{\small\textit{Правильный ответ: Ситуацию, когда разнообразие популяции исчезает слишком рано и поиск замирает.}}

\item Какую задачу решает декодирование хромосомы?
\begin{itemize}[label=$\square$, leftmargin=1.5cm, itemsep=0.2em, topsep=0.3em]
\item Переход от генотипа к фенотипу
\item Уменьшение числа поколений
\item Вычисление евклидова расстояния между кластерами
\item Построение U-матрицы
\end{itemize}
{\small\textit{Правильный ответ: Переход от генотипа к фенотипу.}}

\item Какой из факторов НЕ является обязательным элементом классического ГА?
\begin{itemize}[label=$\square$, leftmargin=1.5cm, itemsep=0.2em, topsep=0.3em]
\item Вычисление градиента целевой функции
\item Популяция решений
\item Функция приспособленности
\item Операторы отбора и вариации
\end{itemize}
{\small\textit{Правильный ответ: Вычисление градиента целевой функции.}}

\end{enumerate}
